%    Latex Template for ACA 2015
%%%%%%%%%%%%%%%%%%%%%%%%%%%%%%%%%%%%%%%%
%    DO NOT CHANGE THE SETTINGS BELOW
%%%%%%%%%%%%%%%%%%%%%%%%%%%%%%%%%%%%%%%%
\documentclass[11pt]{article}
\usepackage{epsfig}
\usepackage{graphicx}
\usepackage[latin1]{inputenc}
\usepackage[T1]{fontenc}
\usepackage{mathptmx}
\usepackage{url}



\begin{document}



%%%%%%%%%%%%%%%%%%%%%%%%
%%%%%%%%%%%%%%%%%%%%%%%%%%%%%%%%%%%%%%%%
%    DO NOT CHANGE THE SETTINGS ABOVE
%%%%%%%%%%%%%%%%%%%%%%%%%%%%%%%%%%%%%%%%
%
\vspace{-0.8cm}
\begin{flushleft}
\Large \textbf{\noindent
%%% PLEASE INSERT THE TITLE OF YOUR TALK HERE %%%%%%%
Computer-Algebra-Aided Chebyshev Methods for Ordinary Differential Equations}
%%%%%%%%%%%%%%%%%%%%%%%%%%%%%%%%%%%%%%%%%%%%%%%%%%%%%%%%
\\
%%%%%%%%%%%%%%%%%%%%%%%%%%%%%%
\vspace{0.5cm}
\normalsize
 %%%%%%             NAMES OF AUTHORS         %%%%%%
 %%%%%% NAME OF SPEAKER SHOULD BE UNDERLINED %%%%%%%%%
\normalsize{
 %%%%%% FOR EXAMPLE %%%%%
 \underline{M. Xue}$^1$
 %%%%%%%%%%%%%%%%%%%%%%%%
} \\
\vspace{5mm}
\textit{\footnotesize
 %%%%%% AFFILIATION OF AUTHORS %%%%%%
$^1$ Vroom Laboratory for Advanced Computing, USA,
 mxue@vroomlab.com\\
 %%%%%%%%%%%%%%%%%%%%%%%%%%%%%%%%%%%%%
}
\end{flushleft}
The solution of ordinary differential equation can be approximated by a
linear combination of so called basis functions. Using the Chebyshev Polynomials
as the basis functions,  the approximation can be expressed as 
\begin{equation}
y(x)=\sum_{r=0}^{\infty} a_{r}T_r(x)
\end{equation}
where $T_{r}(x)$'s are Chebyshev Polynomials of degree $r$, and $a_{r}$'s are the 
coefficients to be determined. In practice, we seek the approximation using a truncated 
expression of (1), namely,
\begin{equation}
\sum_{r=0}^{n} a_{r}T_r(x)
\end{equation}
An online Computer Algebra System (CAS) is used to generate and subsequently solve a system of 
equations concerning a finite number of $a_r$'s. The use of CAS allows the retention of 
more $a_r$'s in (2). It also obviates the need for the traditional pad and pencil 
computations.

Examples will be given to illustrate this approach in solving initial value problems, 
boundary value problems as well as eigenvalue problems for ordinary differential equations
whose coefficients and other terms are themselves polynomials.


%Use this format (without changing the settings) to prepare your  abstract. The abstracts are used for evaluation purposes. Therefore, it is important that you describe  your intended presentation concisely and clearly, including the objectives, status, methodology, results, and significance of your study.
% Abstracts may contain equations, figures and references (see Eqs.~(\ref{eq1}, \ref{eq2}) and Refs.~\cite{KR,WK,ABS}). {\bf The extended abstract length is limited to 5  pages}.  Equations are to be centered horizontally, with equation numbers aligned with the right margin.  Paragraph text is full-justified.  Do not explicitly begin by writing the word ``abstract''.  Equations should be properly punctuated.  For instance, as
%\begin{equation}
%y=3x^2\  \  \  \mbox{and}\  \  \  x>0, \label{eq1}
%\end{equation}
%we see that
%\begin{equation}
%z=\sqrt{y}=\sqrt{3}x. \label{eq2}
%\end{equation}

%\begin{eqnarray}
%\lim_{x \to +\infty} {\rm Abstract}(x) = 1\ {\rm page}.
%\label{eq1}
%\end{eqnarray}

%Your abstract should be prepared in \LaTeX\ according to the template provided  and then \textbf{converted to PDF for submission.}   Abstracts not complying with this template may not be included in the book of abstracts of the ACA 2015. Abstract submission should be done  via the session organizers, they will confirm your submission. If your abstract is accepted, you will be invited to present your talk.
%The book of abstracts will be available to the participants prior to the conference and will provide the audience with more information
%about the papers being presented at the ACA 2015 Workshop. At least one author must be registered for your talk to be included in the program of the conference.

%%%%%%%%%%%%%%%%%%%%%%%%%%%%%%%%%%%%%%%


%%%%%%%%%%%%%%%%%%%%%%%%%%%%%%%%%%%%%%%%
%% BIBLIOGRAPHY
%%%%%%%%%%%%%%%%%%%%%%%%%%%%%%%%%%%%%%%%
\begin{thebibliography}{00}
\addcontentsline{toc}{chapter}{References}
\setlength{\itemsep}{-1mm}
\footnotesize
\bibitem{WK} L. Fox and D.F. Mayers, \textit{Numerical Solution of Ordinary Differential 
Equations}, pp. 179-197 (1987).
\bibitem{ABS} G.H. Golub and J.M. Ortega, \textit{Scientific Computing and Differential Equations}, 
pp. 179-185 (1992).
\bibitem{O}\textit{Omega: A Computer Algebra System Explorer}, at http://www.omega-math.com
\end{thebibliography}
\normalsize
\end{document} 
